\section{From triangle mass to independent sets}
\label{sec:closure}

We now use the triangle estimate to choose independent vertices at a
controlled cost in $F_G^*$. For an $n$-vertex $K_r$-free graph of maximum
degree at most $\Delta\ge3$,
a constant test weight gives $F_G^*$ at least a constant multiple of
$n(\log\Delta)^2/\Delta$. The triangle estimate will provide a vertex
whose closed-neighborhood deletion costs at most a constant times
$\log\Delta$. Induction compares these two bounds and gives the required
independent set. We keep $\Delta$ fixed throughout the induction, then
pass to average degree by deleting high-degree vertices.
Let $B_r\ge0$ be the constant in~\eqref{eq:optimizer-triangle}, and set
\begin{equation}
 D_r=3+\frac32 B_r.
 \label{eq:independence-constant}
\end{equation}

\begin{proposition}\label{prop:maximum-degree}
For every integer $r\ge4$, every real number $\Delta\ge3$, and every finite $K_r$-free graph $G$
with $n$ vertices and maximum degree at most $\Delta$,
\[
 \alpha(G)\ge \frac{F_G^*}{D_r\log\Delta}
 \qquad\text{and}\qquad
 \alpha(G)\ge \frac{n\log\Delta}{4D_r\Delta}.
\]
\end{proposition}

\begin{proof}
Fix $\Delta\ge3$.  We prove the first inequality by induction on $n$.
For the empty graph both sides are zero.  Suppose $n>0$, let $w$ maximize
$F_G$, and write $W=\sum_v w_v$.  By Lemma~\ref{lem:optimizer},
$0<w_v\le1$.

For a vertex $v$, put $A_v=\{v\}\cup N(v)$ and let
$G-A_v$ denote the induced graph on the remaining vertices.
Set $q=0$ on $A_v$ and $q=w$ elsewhere.  The variational identity
for $F_G(w)-F_G(q)$ gives the exact deletion cost
\begin{equation}
 F_G^*-F_{G-A_v}(w|_{V(G)\setminus A_v})
 =w(A_v)+M_{G[A_v]}(w).
 \label{eq:closed-neighborhood-cost}
\end{equation}
Indeed, the vertex divergence is $w(A_v)$, and the quadratic term
$M_G(q-w)$ is $M_{G[A_v]}(w)$.

Average this cost over $v$ with probability $w_v/W$.
An edge $ab$ lies in $G[A_v]$ precisely when $v=a$, $v=b$, or
$v$ is a common neighbor of $a$ and $b$.  Hence
\begin{align}
 \sum_v\frac{w_v}{W}
       \bigl(w(A_v)+M_{G[A_v]}(w)\bigr)
 &=\frac{\sum_v w_v^2+2M_G(w)
       +\sum_v w_v^2w(N(v))+3T_G(w)}{W} \notag\\
 &\le 1+(4+3B_r)\frac{M_G(w)}{W}.
 \label{eq:averaged-deletion-cost}
\end{align}
Here $w_v\le1$ gives $\sum_vw_v^2\le W$ and
$\sum_vw_v^2w(N(v))\le2M_G(w)$; we then used the triangle
bound~\eqref{eq:optimizer-triangle}.

It remains to bound $2M_G(w)/W$ by $\log\Delta$. First, stationarity
and Jensen's inequality for $-\log$ give
\begin{equation}
 \log\frac{n}{W}
 \le \frac1n\sum_v\log\frac1{w_v}
 =\frac1n\sum_v w(N(v))
 =\frac1n\sum_v \deg_G(v)w_v
 \le \frac{\Delta W}{n}.
 \label{eq:optimizer-total-weight}
\end{equation}
If $W/n<1/\Delta$, then the left side exceeds $\log\Delta>1$
and the right side is less than $1$, a contradiction. Therefore
$W\ge n/\Delta$. Applying the concavity of $\log$ with probabilities
$w_v/W$ now gives
\begin{equation}
 \frac{2M_G(w)}{W}
 =\sum_v\frac{w_v}{W}\log\frac1{w_v}
 \le \log\left(\sum_v\frac{w_v}{W}\frac1{w_v}\right)
 =\log\frac{n}{W}
 \le\log\Delta.
 \label{eq:optimizer-edge-weight}
\end{equation}
Thus the average in \eqref{eq:averaged-deletion-cost} is at most
\[
 1+\left(2+\frac32B_r\right)\log\Delta
 \le D_r\log\Delta,
\]
where we used $\log\Delta\ge1$.

Choose a vertex whose cost is at most the average.  Since the restriction
of $w$ is a feasible weight vector for $G-A_v$,
\[
 F_{G-A_v}^*\ge F_G^*-D_r\log\Delta.
\]
The graph $G-A_v$ is $K_r$-free and still has maximum degree at most the
same number $\Delta$.  The induction hypothesis, followed by adjoining
$v$ to an independent set of $G-A_v$, yields
\[
 \alpha(G)\ge 1+\alpha(G-A_v)
 \ge 1+\frac{F_{G-A_v}^*}{D_r\log\Delta}
 \ge \frac{F_G^*}{D_r\log\Delta}.
\]

For the second inequality, test $F_G$ on the constant vector with entry
$p=(\log\Delta)/\Delta$.  Since $|E(G)|\le n\Delta/2$,
\[
 F_G^*\ge np\left(1-\log p-\frac{\Delta p}{2}\right)
 =np\left(1+\frac12\log\Delta-\log\log\Delta\right)
 \ge \frac{np\log\Delta}{4}.
\]
The last inequality holds for every $\Delta>1$: the function
$1+t/4-\log t$ on $t>0$ has minimum $2-\log4>0$ at $t=4$.
Combining this bound with the first inequality proves the result.
\end{proof}

\begin{proof}[Proof of Theorem~\ref{thm:main}]
Let $G$ have $n$ vertices and average degree $d\ge2$.  Delete the vertices
whose degrees in $G$ exceed $2d$, and let $H$ be the induced graph on
the remaining vertices.  Fewer than $n/2$ vertices are deleted, since
the sum of all degrees is $nd$.  Thus $|V(H)|\ge n/2$, and
$H$ is $K_r$-free with maximum degree at most the real number $2d$.
Proposition~\ref{prop:maximum-degree} gives
\[
 \alpha(G)\ge\alpha(H)
 \ge\frac{|V(H)|\log(2d)}{8D_r d}
 \ge\frac{n\log d}{16D_r d}.
\]
We may therefore take $c_r=1/(16D_r)$.
All constants used in obtaining $B_r$, and consequently $D_r$ and $c_r$,
depend only on $r$.
\end{proof}
